% Formato del título de capítulos y secciones
\titleformat{\chapter}[block]{\normalfont\huge\bfseries}{\thechapter.}{.5em}{\Huge}[\vspace{2pt}{\titlerule[2pt]}]

\titlespacing*{\chapter}{0pt}{-19pt}{25pt}

\titleformat{\section}[block]{\normalfont\Large\bfseries}{\thesection.}{.5em}{\Large}

\titleformat{\part}[block]{\titlerule[2pt]\normalfont\Huge\bfseries\centering}{Parte \Roman{part}\\\vspace{15pt}}{0pt}{\Huge}[\vspace{2pt}{\titlerule[2pt]}]

% Tamaños y estilos de elementos en la TOC
\DeclareTOCStyleEntry[
    linefill=\bfseries\TOCLineLeaderFill,
    beforeskip=12pt,
    entrynumberformat=\chapterprefixintoc,
    entryformat=\chaptertocformat,
    pagenumberformat=\chaptertocformat,
    dynnumwidth
]{tocline}{chapter}

\DeclareTOCStyleEntry[
    % linefill=\bfseries\TOCLineLeaderFill,
    beforeskip=30pt,
    entrynumberformat=\chapterprefixintoc,
    entryformat=\parttocformat,
    pagenumberformat=\partpagetocformat,
    numwidth=0pt
]{tocline}{part}

\newcommand\chaptertocformat[1]{\large{\textbf{#1}}}%
\newcommand\chapterprefixintoc[1]{#1}%
\newcommand\parttocformat[1]{\Large{\textbf{#1}}}%
\newcommand\partpagetocformat[1]{} % Don't print the page number for parts

% Alias para estilos de texto comunes
\newcommand{\negritas}[1]{\textbf{#1}}
\newcommand{\cursiva}[1]{\textit{#1}}
\newcommand{\codigo}[1]{\texttt{#1}}

% Formato del código fuente con lstlisting
\lstset{
  basicstyle=\ttfamily,
  breaklines=true,
}

% Márgenes
\geometry{
    a4paper,
    margin=2.75cm
}
\setlength{\marginparwidth}{2cm}

% Limite de profundidad del índice
\setcounter{tocdepth}{2}

% Eliminar el guionado
\tolerance=1
\emergencystretch=\maxdimen
\hyphenpenalty=10000
\hbadness=10000

% Indentación de párrafos
\setlength{\parindent}{.75cm}

\renewcommand{\lstlistingname}{Extracto de código}
\renewcommand*{\lstlistlistingname}{Índice de extractos de código}

\usepackage{newfloat}
\usepackage{caption}

\DeclareFloatingEnvironment[
    fileext=loe,
    listname={Índice de evidencias},
    name=Evidencia,
    placement=H,
    within=chapter
]{evidence}

% Definición de colores personalizados para los enlaces
\definecolor{primaryCyan}{RGB}{85, 204, 170}
\definecolor{primaryMagenta}{RGB}{169, 70, 171}
\definecolor{secondaryCyan}{RGB}{0, 121, 118}
\definecolor{secondaryMagenta}{RGB}{102, 20, 101}

\hypersetup{
    colorlinks=true,       
    linkcolor=primaryMagenta,   % Color para referencias internas 
    citecolor=secondaryCyan,   % Color para citas bibliográficas 
    urlcolor=primaryCyan      % Color para direcciones web 
}

% Comandos para establecer variables
\newcommand{\setTitle}[1]{\def\tfgTitle {#1}}
\newcommand{\setAuthor}[1]{\def\tfgAuthors {#1}}
\newcommand{\setDegree}[1]{\def\tfgDegree {#1}}
\newcommand{\setSupervisor}[1]{\def\tfgSupervisor {#1}}
\newcommand{\setDepartment}[1]{\def\tfgDepartment {#1}}
\newcommand{\setMonth}[1]{\def\tfgMonth {#1}}
\newcommand{\setYear}[1]{\def\tfgYear {#1}}
\newcommand{\setDedication}[1]{\def\tfgDedication {#1}}

\newcommand{\covRed}[1]{\textcolor{red!80!black}{\textbf{#1}}}
\newcommand{\covCross}{\textcolor{red!80!black}{\scalebox{1.4}{\textbf{\texttimes}}}}
\definecolor{covYellow}{RGB}{218,165,32}
\newcommand{\covYellow}[1]{\textcolor{covYellow}{\textbf{#1}}} 
\newcommand{\covPartial}{\textcolor{covYellow}{\raisebox{0.25ex}{\scalebox{1.8}{\textbf{\texttildelow}}}}}
\newcommand{\covGreen}[1]{\textcolor{green!60!black}{\textbf{#1}}}
% \newcommand{\covCheck}{\textcolor{green!60!black}{\textbf{\checkmark}}}
\newcommand{\covCheck}{\textcolor{green!60!black}{\ding{51}}}       
\definecolor{bashYellow}{RGB}{204, 204, 0} 
\newcommand{\bashYellow}[1]{\textcolor{bashYellow}{\textbf{#1}}}
\definecolor{bashGrey}{RGB}{140, 140, 140}
\newcommand{\bashGrey}[1]{\textcolor{bashGrey}{\textbf{#1}}}
\newcommand{\cliPrompt}{\textcolor{bashGrey}{\textbf{>}}}


\lstdefinestyle{terminal}{
    language={},
    basicstyle=\fontsize{12pt}{13.8pt}\ttfamily\color{black},
    keywordstyle=\color{black},
    commentstyle=\color{black},
    stringstyle=\color{black},
    identifierstyle=\color{black},
    numbers=none,
    frame=none,
    columns=fullflexible,
    keepspaces=true,
    showstringspaces=false,
    escapeinside={(*@}{@*)}
}
\lstset{style=terminal}

\makeatletter
\renewcommand{\@pnumwidth}{2.2em}
\renewcommand{\@tocrmarg}{2.5em}
\makeatother